\section{Conclusion}\label{sec:conclusion}

LLM agents for process operations need not pay for a full simulation of every proposal, because the forward pass
that produces a proposal already estimates how likely it is to fail. Across two domains, five open-weight models
and a pre-registered protocol with sealed test sets, a screen that reads the agent's internals, with no model of
the plant, removed 11--79\% of validator calls in graph path planning, with certified accepts for four models,
and 12--53\% in CSTR fault recovery, almost entirely through rejections that were 95--98\% precise. Combined with
observable features, internals added up to 15 points of skipped calls where the agent's errors were failures to
consult the relevant information. They added little where the agent consulted the right information and
reasoned incorrectly, a distinction that region-based attention grounding measures directly. In closed loop,
internals-based screens accepted failing retries, a distribution shift outside the certificate; allowing
unchecked execution only for first proposals restored the certificate's premise and kept unsafe executions at
0--2 per 400 episodes while 29--65\% of validator calls were avoided. Because the screen's cost is fixed and the
simulator's grows with fidelity and horizon, these avoided calls translated into compute savings of 21--62\% at
60\,s per validation for four of five models.

Future work should train and certify screens on retries rather than excluding them, certify CSTR screens on larger
partitions, test larger and API-served models through whatever internal signals they expose, and evaluate the
approach on plants with richer fault libraries, real-time constraints and model mismatch between the digital twin
and the plant.
